\subsection{Milestone 4 -- Refactoring automatizzato}

La Milestone 4 caratterizza l'efficacia del refactoring automatizzato via LLM su classi
Java di OpenJPA, articolata in selezione delle classi target (Step~8) e produzione
iterativa di versioni refactored con test progressivamente più stringenti (Step~9,
Tabella~\ref{tab:m4-versions}).

\subsubsection{Selezione delle classi (Step 8)}

Tutte le classi Java di produzione dell'ultima release disponibile nel repository
Git (\texttt{4.1.1}) vengono analizzate con PMD 7.7.0 (stessi ruleset di M1) per
contare le violazioni (\emph{Nsmells}). Vengono escluse le classi con LOC $< 100$ (scarsità di metodi), le interfacce Java
e le classi astratte (entrambe prive di implementazione completa: un LLM non può
applicare refactoring significativo su un tipo non istanziabile).
Le restanti classi vengono ordinate per Nsmells decrescente; a parità di Nsmells,
criterio secondario LOC decrescente. La selezione segue la formula del docente:
$G \Rightarrow 7 \bmod 5 = 2$, caso~2 $\Rightarrow$ rank~$3$ e rank~$N{-}2$
(Tabella~\ref{tab:m4-selected}).

\subsubsection{Refactoring automatizzato (Step 9)}

Per ogni classe si produce la sequenza $C_0 \!\to\! C_1 \!\to\! C_2 \!\to\! C_3
\!\to\! C_4$ (Tabella~\ref{tab:m4-versions}) fornendo all'LLM il codice originale,
le diagnostiche SonarCloud e i test scritti dallo studente in modo incrementale.
Il prompt impone: rimozione degli smell indicati senza alterare la funzionalità,
superamento dei test forniti, nessuna modifica aggiuntiva, compatibilità con il sistema.
